\subsubsection{How OPALX calculates the one-turn matrix}
\label{sec:sc-one-turn-matrix}
The one-turn matrix answers a local question: if a particle starts a little
away from the closed orbit, how does that small displacement change after
one revolution? OPALX computes this matrix by differentiating a numerically
tracked return map. It uses the full nonlinear external magnetic fields for
each trajectory, then extracts their \emph{linear} response near the closed
orbit. The native \texttt{COF} calculation shown here has no space charge.
It provides the zero-current reference for the KV and particle comparisons.

\paragraph{Define the observation section and coordinates.}
A return map must start and end at the same observation section. Here that
section is the plane $Z=0$, crossed in the forward direction. The section's
local axes coincide with the global axes at launch: $x_{\rm sec}=X$,
$y_{\rm sec}=Y$, and its forward axis is $Z$. In particular, its horizontal
coordinate is radial at launch; it differs from the moving orbit-normal
coordinate used in the KV kicks. OPALX stores the transverse state as
\begin{equation}
 \vect u=\begin{pmatrix}x_{\rm sec}\\\pi_{x,\rm sec}\\
                        y_{\rm sec}\\\pi_{y,\rm sec}\end{pmatrix},
 \qquad \pi_{x,\rm sec}=\frac{p_{x,\rm sec}}{mc},\qquad
        \pi_{y,\rm sec}=\frac{p_{y,\rm sec}}{mc}.
 \label{eq:cof-coordinates}
\end{equation}
Positions are in metres; the two normalized \emph{mechanical momenta} are
dimensionless. They are not the slopes $x'$ and $y'$. At fixed kinetic energy,
the total normalized momentum is $P_\star=\beta\gamma=0.16044471664$ and each
perturbed launch uses
\begin{equation}
 \pi_{z,\rm sec}=+\sqrt{P_\star^2-\pi_{x,\rm sec}^2
                                      -\pi_{y,\rm sec}^2}.
 \label{eq:cof-longitudinal-momentum}
\end{equation}
Thus changing a transverse momentum does not change the launch energy.
This construction gives a four-dimensional transverse map at fixed energy,
rather than a six-dimensional map including arrival time and energy spread.

\paragraph{Track the return map and find its fixed point.}
Write the result of one directed return as $\mathcal F(\vect u)$.
OPALX launches a ray from the section, tracks it through the ring, and detects
the next forward return after the ray has left the launch neighbourhood.
It brackets the section crossing and refines its time by bisection, advancing
the ray again to each trial time. Every perturbed ray therefore ends on the
same geometric section, even if its flight time differs slightly. Ending
all rays after a common, approximate revolution time would introduce an
arrival-time displacement into the transverse map.

The closed orbit is the fixed point $\vect u_\star$ satisfying
\begin{equation}
 \mathcal F(\vect u_\star)=\vect u_\star.
 \label{eq:cof-fixed-point}
\end{equation}
The \texttt{COF} command solves this condition with Newton corrections and a
line search. At a trial state $\vect u$, the correction obeys
\begin{equation}
 (M-\mathbb I)\,\delta\vect u
       =-[\mathcal F(\vect u)-\vect u],
 \label{eq:cof-newton}
\end{equation}
where $M$ is the derivative of $\mathcal F$ at that trial state. OPALX scales
the coordinates when solving this system, reduces the correction when
necessary, and verifies convergence with a fresh return-map evaluation.
For this machine, the converged reference state is
\begin{equation}
 \vect u_\star=\begin{pmatrix}
 -0.0270744445304\ \mathrm m\\
  0.0162803888787\\0\ \mathrm m\\0
 \end{pmatrix},\qquad
 \pi_{z,\star}=0.159616590728.
 \label{eq:cof-reference-state}
\end{equation}
Its radius is $R_0+x_{\rm sec,\star}=4.39292555547$ m. The nonzero horizontal
momentum shows that the orbit is not tangent to the placement circle at the
launch point.

\paragraph{Differentiate with paired rays.}
For a small perturbation $\delta\vect u$ about the closed orbit,
\begin{equation}
 \mathcal F(\vect u_\star+\delta\vect u)-\vect u_\star
       =M\,\delta\vect u+\mathcal O(\|\delta\vect u\|^2),\qquad
 M_{ij}=\left.\frac{\partial\mathcal F_i}{\partial u_j}
                                                   \right|_{\vect u_\star}.
 \label{eq:cof-linearization}
\end{equation}
OPALX perturbs one coordinate at a time, leaving the other three fixed, and
tracks both signs of that perturbation. Column $j$ follows from the centred
finite difference
\begin{equation}
 \boxed{M_{ij}\simeq
 \frac{\mathcal F_i(\vect u_\star+h_j\vect e_j)
             -\mathcal F_i(\vect u_\star-h_j\vect e_j)}{2h_j}.}
 \label{eq:cof-finite-difference}
\end{equation}
The vector $\vect e_j$ selects coordinate $j$. Four columns require eight
perturbed return trajectories. The subtraction removes the common reference
motion; dividing by $2h_j$ measures the response per unit launch perturbation.
The \texttt{FDSTEP} values set $h_j$. A large perturbation includes nonlinear
motion, whereas a very small one amplifies tracking and roundoff errors.
Both the tracking step and the perturbations need convergence checks.

\paragraph{The matrix for the 16-cell, 12 MeV machine.}
The ray integrator is DOP853. The refined native calculation uses
\texttt{DT}=4 ns, position perturbations of $5\ \mu$m, and normalized-momentum
perturbations $5\times10^{-6}P_\star=8.0222\times10^{-7}$.
The resulting matrix, in exactly the coordinate order of
Eq.~\eqref{eq:cof-coordinates}, is
\begin{equation}
 \boxed{M\simeq\begin{pmatrix}
 -0.81100424 & 3.14919254 & 0 & 0\\
 -0.09038064 &-0.88208415 & 0 & 0\\
 0 & 0 &-0.92829272 & 6.05362895\\
 0 & 0 &-0.06987036 &-0.62160435
 \end{pmatrix}.}
 \label{eq:cof-machine-matrix}
\end{equation}
Within each $2\times2$ block, the diagonal entries are dimensionless, the
upper-right entry has units m, and the lower-left entry has units m$^{-1}$.
For example, a 1 mm horizontal position perturbation with no initial change
in transverse momentum returns, to first order, as a $-0.811$ mm position
perturbation and a $-9.04\times10^{-5}$ change in $p_x/(mc)$.
The zero off-diagonal blocks express horizontal--vertical decoupling of the
linear motion about this midplane orbit. They do not imply that the nonlinear
motion at finite amplitude is uncoupled.

The determinants of the two blocks, computed before rounding the displayed
entries, are $\det M_x=0.9999999990$ and $\det M_y=1.0000000044$.
Repeating the calculation with 8 ns tracking steps and twice the
finite-difference perturbations gives a largest change of
$5.8\times10^{-7}$ m in the vertical upper-right matrix element; the
horizontal and vertical tune changes are
$5.2\times10^{-11}$ and $8.8\times10^{-9}$. This joint refinement supports
the zero-current reference but does not separate the two error sources or
establish every displayed matrix digit independently. DOP853 is the ray
integrator selected for this \texttt{COF} calculation; the collective beam
runs use the Boris integrator described below.

\paragraph{Extract the tunes and recover the missing integer.}
For each stable, approximately area-preserving plane block,
\begin{equation}
 \lambda_{j,\pm}\simeq e^{\pm i\mu_{j,1}},\qquad
 \cos\mu_{j,1}\simeq\frac{\operatorname{Tr}M_j}{2},\qquad
 q_j=\frac{\mu_{j,1}}{2\pi}.
 \label{eq:cof-eigenphase}
\end{equation}
Here $\mu_{j,1}$ is the principal one-turn eigenphase. In the numerical
extraction OPALX normalizes by $\sqrt{\det M_j}$ when determining the cosine
and uses the sign of the upper-right entry to orient the phase. For this
machine the fractional tunes are $q_x=0.41066045$ and $q_y=0.39111213$.
The trace alone gives a cosine, with a sign ambiguity for the phase, and
eigenvalues cannot distinguish a phase from that phase plus an integer
multiple of $2\pi$.

To recover the full tune, OPALX also follows the reference ray and paired
shadow rays through the ring, projects their separation onto the local
horizontal normal and the vertical direction, and unwraps the associated
linear phase. It checks the phase sampling and consistency with the endpoint
matrix. The winding supplies the integer part; the accurate endpoint matrix
supplies the fractional part. The result is
\begin{equation}
 \boxed{Q_x=3.41066045,\qquad Q_y=3.39111213.}
 \label{eq:cof-full-tunes}
\end{equation}
This full-tune reconstruction assumes uncoupled planar motion. It adds
information that is absent from the one-turn eigenvalues alone. These are
zero-current, infinitesimal-amplitude tunes; the finite-amplitude NAFF
footprints and the intensity-dependent KV tunes address the additional
physics discussed in the following subsections.
